% =====================================================================
% Seção de resultados.
%
% As tabelas abaixo são GERADAS por report.py a partir dos JSON de
% results/validation/. Não editar os .tex de tabela à mão: rode
%     python validation.py --suite all
%     python report.py
% e copie results/validation/paper/*.tex e *.pdf para junto do artigo.
%
% Números no texto conferidos contra a execução de 2026-10-15,
% commit d57f4d0, PyTorch 2.6.0+cu118, RTX 3060.
% =====================================================================

\section{Results}
\label{sec:results}

Before the individual findings, one property of the design bounds how all of
them may be read. With five folds, the smallest two-sided $p$ a Wilcoxon
signed-rank test can return is $0.062$; no comparison in this section can
therefore reach the conventional $0.05$ threshold, whatever the size of the
underlying effect. This is a property of the sample of folds, not of the
effects, and it is why we report Cohen's $d$ and the number of folds that moved
in each direction as the primary evidence, with $p$ included only for
completeness.

\subsection{Generality across datasets}

Table~\ref{tab:generality} applies the identical core and protocol to four
datasets: two natural-image benchmarks~\cite{lecun1998mnist,xiao2017fashion} and two medical ones
drawn from MedMNIST~v2~\cite{yang2023medmnist}: BloodMNIST, from peripheral blood
smears~\cite{acevedo2020blood}, and DermaMNIST, from dermatoscopic
images~\cite{tschandl2018ham10000}.

\input{table_generality}

A single configuration, chosen once and never retuned per dataset, reaches
$0.986$ accuracy on MNIST, $0.895$ on Fashion-MNIST and $0.937$ on BloodMNIST.
The last of these sits about two points below the published ResNet-18 reference
for the same data at the same resolution, which is a reasonable showing for an
architecture that routes everything through eight discrete symbols on the way
to its decision.

DermaMNIST is the exception. Accuracy reaches $0.726$, close to the
$0.735$ of the ResNet-18 reference, while balanced accuracy is $0.469$ and the
pair gain is slightly \emph{negative} ($-0.010$). The dataset has an imbalance
ratio of $58{:}1$, and the model is substantially predicting the majority class;
accuracy alone would have reported this as a near-baseline success.
Figure~\ref{fig:generality} shows the divergence directly. We keep this case
because it is the reason the protocol reports balanced accuracy and pair gain
alongside accuracy: an evaluation without them would have scored a collapsed
model as a working one.

\begin{figure}[!t]
\centering
\includegraphics[width=0.99\linewidth]{fig_generality.pdf}
\caption{Accuracy and balanced accuracy under the same core and protocol.
Dashed lines mark the published ResNet-18 accuracy for the medical datasets.
The gap between the two bars on DermaMNIST is majority-class collapse: the
dataset is imbalanced $58{:}1$.}
\label{fig:generality}
\end{figure}

\subsection{Encoder and bottleneck}

Table~\ref{tab:encoder} separates two things that are easy to confuse: the
quality of the features entering the core, and the price of the discrete
bottleneck itself. The third variant per dataset removes the core entirely,
keeping the same encoder and the same program dimensionality, so the difference
isolates discretisation.

\input{table_encoder}

The encoder effect is large and consistent. Replacing the dense layer over
flattened pixels with a small CNN moves Fashion-MNIST from $0.858$ to $0.895$
($d=4.94$, 5/5 folds) and BloodMNIST from $0.860$ to $0.937$ ($d=10.68$, 5/5
folds). This is worth stating plainly because an earlier version of these
experiments used the dense encoder throughout and therefore understated the
architecture: what looked like a limitation of the sparse VSA core was mostly a
limitation of what we were feeding it.

The bottleneck effect is smaller. Against the no-core
control, \emph{accuracy is statistically indistinguishable}: the control is
ahead by $0.004$ on Fashion-MNIST ($d=0.61$, 3/5 folds) and behind by $0.002$ on
BloodMNIST ($d=-0.24$, 2/5 folds). Compressing the representation into eight
groups of sixteen symbols costs essentially nothing in decision quality, in
either direction, on either dataset.

The cost appears instead in \emph{ranking}. Macro AUC favours the control
consistently, by $0.008$ on both datasets, with large effect size and all five
folds agreeing in each case ($d=9.05$ and $d=4.99$). The interpretation we
offer is that discretisation preserves the decision boundary while discarding
the graded confidence around it: an argmax over sixteen symbols per group cannot
express the fine differences in certainty that a dense vector carries, and AUC
is the metric that reads those differences. For a downstream use that
needs calibrated scores this matters; for one that needs decisions, it does not.

\subsection{Redundancy penalty}

Trained only to classify, the groups converge on redundant codes: inter-group
NMI sits at $0.54$ on Fashion-MNIST and $0.53$ on BloodMNIST, with pair gains of
only $0.068$ and $0.053$. In that regime the pairwise half of the program adds
little to the unary half. Table~\ref{tab:specialization} and Figure~\ref{fig:tradeoff}
trace what the redundancy penalty of Equation~\ref{eq:redundancy} changes.

\input{table_specialization}

\begin{figure}[!t]
\centering
\includegraphics[width=0.99\linewidth]{fig_tradeoff.pdf}
\caption{Effect of the redundancy penalty. The two datasets track each other
closely in all three panels. Compositionality (left) rises sharply and then
plateaus; redundancy (centre) falls monotonically; readability of the
hypervector (right) degrades and does not recover at larger weights.}
\label{fig:tradeoff}
\end{figure}

\textbf{Effect of the penalty.} At $\lambda=0.25$ the
pair gain more than doubles on both datasets, from $0.068$ to $0.159$ on
Fashion-MNIST ($d=7.87$, 5/5 folds) and from $0.053$ to $0.133$ on BloodMNIST
($d=9.39$, 5/5 folds), while inter-group redundancy halves. Accuracy moves by
$+0.002$ and $-0.003$ respectively, well inside the fold-to-fold spread. The two curves in
Figure~\ref{fig:tradeoff} nearly coincide across a greyscale clothing dataset
and an RGB microscopy dataset, which is the replication evidence available at
this stage.

\textbf{Larger weights.} The pair gain peaks at $\lambda=0.25$ on
both datasets and then flattens or declines, even as redundancy keeps falling
to $0.20$ and $0.18$ at $\lambda=1$. Pushing codes towards independence
eventually removes the shared structure the objective was meant to create:
perfectly independent groups have nothing to compose. The useful range of $\lambda$ is
narrow and at the low end.

\textbf{Effect on readability.} VSA recovery falls from $0.988$ to
$0.872$ on Fashion-MNIST and from $0.991$ to $0.780$ on BloodMNIST at
$\lambda=0.25$, and continues to degrade with larger weights. Code perplexity
falls alongside it, from $6.07$ to $4.39$ and from $4.90$ to $3.66$, which is
consistent with independence being bought by relaxing the separation between
codebook entries, since cleanup (Equation~\ref{eq:cleanup}) is a cosine
lookup against those same entries.

The severity of this cost is sensitive to the training budget, and we report
this because an earlier run misled us. Under a shorter schedule and the dense
encoder, recovery collapsed to roughly $0.65$ at every non-zero weight, which
suggested a hard discontinuity. With the CNN encoder and four times the
optimisation steps, unpenalised recovery rises to $0.99$ and the penalised
value at $\lambda=0.25$ retains $0.87$ on Fashion-MNIST. The conflict between
independence and readability is real and reproducible, but it is a graded
trade-off, not the cliff the earlier evidence implied.

This distinction matters for the architecture's purpose. Composing two
independently trained networks requires reading codes back out of a
superposition; a program that classifies well but decodes poorly is of limited
use for that. At $\lambda=0.25$ the construction buys a doubled pair gain for
roughly twelve points of recovery on Fashion-MNIST and twenty-one on
BloodMNIST, and whether that is a good trade cannot be settled until the
composition step itself is measured.

\subsection{Statistical summary}

Table~\ref{tab:stats} collects the fold-paired comparisons behind the claims
above. Every configuration of a dataset shares its fold partitions, so the
differences are paired by construction.

\input{table_stats}

\subsection{Limitations}

Four limitations bound what the above establishes.

The composition mechanism itself is untested. The routine that \emph{applies} a
Glue, translating codes between algebras, merging with target-private symbols
and composing programs, is implemented and structurally validated,
but the routine that \emph{estimates} the alignment from data is not yet part of
the reproducible pipeline. Every result here therefore concerns a single
network, and the central claim of the project remains open.

The evaluation is confined to $28\times28$ images and to one core
configuration. The generality claim is about domains, not scale: nothing here
shows the construction survives at higher resolution, with more groups, or on
non-visual input.

Five folds give little statistical power, as noted at the outset. The effects we
rely on are large ($d > 4$ throughout), but the design cannot distinguish a
small real effect from none.

Finally, $\lambda$ was held fixed throughout training. A schedule that anneals
it, or an explicit codebook-separation term, might recover the lost readability;
both are untested and are the natural next experiments.
